% Generated draft: requires mathematical and source review.
\documentclass[11pt]{article}
\usepackage[margin=25mm]{geometry}
\usepackage{fontspec}
\setmainfont{Latin Modern Roman}
\usepackage{amsmath,amssymb,mathtools,unicode-math}
\setmathfont{Latin Modern Math}
\usepackage{enumitem,needspace,xurl,hyperref}
\setlength{\parindent}{0pt}
\setlength{\parskip}{5pt}
\setlength{\emergencystretch}{3em}
\newcommand{\N}{\mathbb{N}}
\newcommand{\Z}{\mathbb{Z}}
\newcommand{\Q}{\mathbb{Q}}
\newcommand{\R}{\mathbb{R}}
\newcommand{\C}{\mathbb{C}}
\newcommand{\F}{\mathbb{F}}
\newcommand{\A}{\mathbb{A}}
\newcommand{\PP}{\mathbb{P}}
\newcommand{\E}{\mathbb{E}}
\newcommand{\eps}{\varepsilon}
\newcommand{\dd}{\,\mathrm{d}}
\begin{document}
\section*{UnsolvedMath ID: 30000933}
\textbf{Source problem code:} catalogue entry
\subsection{Definitions and mathematical statement}
Let \(\Omega\subset\R^d\) be a domain, let \(u:\Omega\to\R^d\) be a velocity field and \(p:\Omega\to\R\) a pressure, and consider the Stokes eigenvalue equations
\[
-\Delta u+\nabla p=\lambda u,\qquad \operatorname{div}u=0\quad\hbox{in }\Omega,
\qquad u=0\quad\hbox{on }\partial\Omega.
\]
The source asks the following related questions.
\begin{enumerate}
\item If \(p\) is constant on the relevant portion \(\Gamma\subseteq\partial\Omega\), under what intended domain and boundary hypotheses does it follow that \(u\equiv0\) and \(p\) is constant throughout \(\Omega\)? The record does not specify \(\Gamma\), the domain restrictions, or the trace and solution regularity needed for this implication.
\item Is the disk the only domain admitting a solution with \(u\not\equiv0\) and \(p\equiv0\)? The mention of a disk indicates a planar classification question, but the source does not explicitly fix the dimension, boundedness, connectedness of the boundary, or boundary regularity. For any fixed solution the equation depends only on \(\nabla p\), so a spatially constant pressure can be normalized to zero.
\item Establish the related unique-continuation assertion for nonsmooth boundaries. The exact boundary class, observation set, hypotheses, and conclusion of this additional assertion are not supplied.
\end{enumerate}
The first question is retained as a request to identify the intended rigidity statement, rather than an unconditional implication for every domain: the supplied assessment explicitly mentions a disk counterexample and known corner results. Their precise statements are absent, so neither the exception class nor a corner theorem is inferred here.
\subsection{Short English statement}
The problem concerns rigidity of Stokes eigenfunctions whose pressure is constant on part of the boundary, the proposed uniqueness of the disk among domains carrying a nonzero zero-pressure eigenfunction, and a nonsmooth-boundary extension. The missing geometric and regularity hypotheses prevent a fully quantified theorem statement.
\subsection{Sources}
\begin{itemize}
\item Local UnsolvedMath registry, permanent ID \(30000933\), source code \texttt{catalogue entry}, ``Stokes Eigenfunctions with Constant Boundary Pressure.'' Registry curation: Ulam AI / UnsolvedMath Contributors, CC BY 4.0; the underlying report and literature retain their supplied provenance.
\item \emph{Optimal Control of Coupled Systems of PDE} (2008), Oberwolfach Reports, pp.\ 587--668, DOI \url{https://doi.org/10.4171/owr/2008/13}.
\item Supplied ``Stokes--Schiffer analysis'': \url{https://math.ucr.edu/~kelliher/DeskDrawer/SchifferAndFluids.pdf}.
\item The supplied literature triage, checked 2026-08-21, labels the problem open, describes its relation to the Schiffer conjecture, says that disk counterexamples and corner results are known, and says the domain characterization remains unresolved in general. These are dated source claims, not an independent assessment.
\end{itemize}
\end{document}
